%\section{问题重述}
%
%\subsection{问题背景}
%
%数学建模是解决实际问题的重要工具。本文研究的问题涉及多维度数据分析与优化决策。
%
%\subsection{问题分析}
%
%\textbf{问题一：}对给定数据集进行分析和预测，要求模型具有泛化能力。
%
%\textbf{问题二：}考虑多因素影响，改进和优化模型，提高预测精度。
%
%\textbf{问题三：}结合实际约束条件，设计最优策略方案。
%
%\textbf{问题四：}结合实际约束条件，设计最优策略方案。
\section{问题重述}

\subsection{问题背景}

人工智能应用扩张使数据中心算力需求与能耗同步增长。数据中心能耗建模需同时考虑IT负荷、设施能效和电力交互\cite{ref1,ref2}，综合能源系统还需协调数据流与能量流\cite{ref3}。各区域的算力、电价、碳强度、新能源和网络时延不同，加之新能源波动及任务时限差异，形成跨区域、跨时段的算--储--电协同优化问题。


\subsection{问题分析}

\textbf{问题一：}预测第2376--2399小时各区域、各类任务的GPU需求；再以实际到达任务为对象，在时限、时延、GPU和功率约束下确定执行区域与开工时刻，形成基础算力基线。

\textbf{问题二：}在实际任务和能源参数下重新安排任务区域与开工时刻，建立兼顾运行成本、碳排放、网络时延和新能源利用率的碳感知调度模型。

\textbf{问题三：}固定题目给定的算力负荷，仅优化新能源分配、储能充放电和购售电，评价储能对成本、碳排放、峰值净购电及波动的影响。

\textbf{问题四：}联合优化任务调度、储能和购售电，在算力、能源与服务质量约束下比较联合方案、顺序基准及不同运行情景。


\subsection{整体技术路线}
如图\ref{fig:framework}，问题一形成纯算力基线，问题二引入电价、碳强度和新能源，问题三在固定负荷下考察储能，问题四联合开放任务与能源决策，构成“基础算力--算电协同--储能调节--算储电联合调度”的递进链路。


\begin{figure}[H]
	\centering
	\includegraphics[width=0.80\textwidth]{figures/华数杯流程图.pdf}
	\caption{多区域算--储--电协同优化整体技术路线}
	\label{fig:framework}
\end{figure}


